\documentclass[11pt, a4paper]{article}
% ---- Core Geometry & Margins ----
\usepackage[margin=1in, top=1.2in, bottom=1.2in]{geometry}
\setlength{\parindent}{0pt}
\setlength{\parskip}{0.8em}
% ---- Font & Micro-Typography ----
\usepackage[utf8]{inputenc}
\usepackage[T1]{fontenc}
\usepackage{lmodern}
% \usepackage{newtxtext}
% \usepackage{newtxmath}
% \usepackage{libertinus}
\usepackage{microtype}

% ---- Language & Math ----
\usepackage[english]{babel}
\usepackage{amsmath, amssymb, amsthm}
\usepackage{booktabs}  % In your preamble

% ---- Formatting & Aesthetics ----
\usepackage{booktabs}
\usepackage{graphicx}
\usepackage{xcolor}
\usepackage{csquotes}

% ---- Color Palette ----
\definecolor{navyblue}{RGB}{16, 55, 87}
\definecolor{darkslate}{RGB}{53, 64, 83}

% ---- Hyperlinks Setup ----
\usepackage[
  pdfencoding=auto,
  colorlinks=true,
  linkcolor=navyblue,
  citecolor=navyblue,
  urlcolor=navyblue
]{hyperref}

% ---- Bibliography Setup ----
\usepackage[
  backend=biber,
  style=numeric-comp,
  sorting=nyt,
  giveninits=true
]{biblatex}
\addbibresource{references.bib}

% --- Metadata ---
\title{Random Thoughts}
\author{Guetchuin Tadouanla}
\date{today}

% ====================================================================================
\begin{document}
\pagenumbering{roman}
\setcounter{page}{1}

\tableofcontents
\newpage
\pagenumbering{arabic}
\setcounter{page}{1}

\section{Introduction}
In terms of sustainable development, Cameroon is currently ranked $\mathbf{135^{th}}$ out of
\textbf{169} countries, with a sustainable development Goal Index score of $\mathbf{57.6\%}$
in accordance to the sustainable development report of 2025. The figure below illustrates the
current situation of Cameroon in regards to the $\mathbf{17}$ Sustainable Development Goals.

% \includegraphics[width=1\textwidth]{figures/SDG_Dashboard_and_Trends.png}
% \includegraphics[scale=0.3]{figures/SDG_dashboard.png}
% \includegraphics[width=17cm, height=6cm]{figures/SDG_Dashboard_and_Trends.png}

\begin{figure}[htbp]
  \centering
  \includegraphics[width=1\textwidth]{figures/SDG_dashboard.png}
  \caption{Cameroon SDG Dashboard and Trends}
  \label{fig:SDG} \end{figure}

The country has achieved two sustainable development goals namely: SDG 12: \textbf{Responsible
  Consumption and Production } and SDG 13:  \textbf{Climate Action}. SDG 14 shows some challanges
remaining while SDG 8 shows significant challenges remaining in order to achieve their sustainable
development goal. The remaining SDGs are facing Major challenges and very far from the sustainable
development goal.


My interest will be on SDGs 7(\textbf{Affordable and Clean Energy}), 9(\textbf{Industry Innovation
  And Infrastructure}) and 11(\textbf{Sustainable Cities And Communities}) as this is where I'm
able to contribute solution ideas to the challenges.

\section{SDG 7: Affordable and Clean Energy}
This sustainable development goal  aims at ensuring access to affordable, reliable, sustainable
and modern enrgy for all. As earlier mentioned, Cameroon is facing Major challenges regarding
the affordability and cleanliness  of the energy consumed in the country. The sustainability
development report shows a stagnation or increasing at a rate of less than $\textbf{50\%}$ of the required
rate towards the goal. This is proven through the stagnating and low trending indicators below.

\begin{figure}[htbp]
  \centering
  \includegraphics[width=1.1\textwidth]{figures/sustainable-development-report-2026-population-with-access-to-electricity.png}
  \caption{Population with access to electricity}
  \label{fig:SDG}
\end{figure}

\begin{figure}[htbp]
  \centering
  \includegraphics[width=1\textwidth]{figures/sustainable-development-report-2026-population-with-access-to-clean-fuels-and-technology-for-cooking.png}
  \caption{Population with access to clean fuels and technology for cooking}
  \label{fig:SDG}
\end{figure}

\begin{figure}[htbp]
  \centering
  \includegraphics[width=1\textwidth]{figures/sustainable-development-report-2026-co2-emissions-from-fuel-combustion-per-total-electricity-output.png}
  \caption{$CO_2$ emissions from fuel combustion per total electricity output}
  \label{fig:SDG}
\end{figure}

\begin{figure}[htbp]
  \centering
  \includegraphics[width=1\textwidth]{figures/sustainable-development-report-2026-renewable-energy-share-in-total-final-energy-consumption.png}
  \caption{Renewable energy share in total final energy consumption}
  \label{fig:SDG}
\end{figure}

\newpage
\subsection{Indicator: Population with access to clean fuels and technology for cooking}
From the above indicators, having a closer look into this indicator reveals connections to other
SDGs.

The hazardous cooking conditions in which low income household in Cameroon cook in, mainly
due to the use of charcoal as fire source, which pollutes the air in the environment in which they
cook in. A possible solution idea to better the condition is through the use of \textbf{LPG}
\textbf{Liquified Petroleum Gas}. The aim is to ease the access to LPG and making it a cheaper
or perhaps equally affordable option for such households in comparism to their current . Which in return

The company \textbf{PayGo Energy} has found a solution to easing the access of LPGs in low income households in Kenya.
The cooking conditions in low income household in Kenya is pretty similar to the ones in Cameroon. So extending this
solution idea or the company itself is a solution for the hazardous cooking conditions.

Another solution idea would be through the collection of biomass and using that as fuel.
Given that some regions of the country mainly the northern part of the country practice
a lot of animal raring from cows to sheeps and goats from which biomass could be sourced from.

\section{SDG 9: Industry Innovation and Infrastructure}
SDG 9 is about Building resilient infrastructure, promote inclusive and sustainable
indsutrialization and foster innovation. This is another challenge Cameroon is currently facing as
they are stagnating or increasing at less than $\mathbf{50\%}$ of the required rate.
There are several stagnating and low trending indicators proving this to be a challenge to the
country's sustainable development.

\begin{itemize}
  \item Rural population with access to all-season roads
  \item Population using the internet
  \item Mobile broadband subscriptions
  \item Logistics Performance Index: Infrastructure Score
  \item Score of top 3 universities in Times Higher Education
  \item Articles published in academic journals
  \item Expenditure on research and development
  \item Total patent applications by applicant's origin
  \item E-government: Online Services Index
\end{itemize}

\subsection{Indicator: Total patent applications by applicant's origin}
The focus point would be directed at this indicator.
\end{document}
